\documentclass[12pt]{exam}
\usepackage[utf8]{inputenc}

\usepackage{amsmath,amstext,amsthm,amssymb,amsxtra}
\usepackage[top=1.5in, bottom=1.5in, left=1.25in, right=1.25in]	{geometry}
%\usepackage[normalem]{ulem}
\usepackage{txfonts} % pxfonts txfonts 
\usepackage[T1]{fontenc}
\usepackage{lmodern}
\renewcommand*\familydefault{\sfdefault}
 \usepackage{euler}   % better than the option below
\usepackage{pdfsync}
\usepackage{multicol}
\newcommand{\ci}[1]{_{ {}_{\scriptstyle #1}}}
\usepackage{graphicx}
\graphicspath{ {images/} }


\newcommand{\norm}[1]{\ensuremath{\left\|#1\right\|}}
\newcommand{\abs}[1]{\ensuremath{\left\vert#1\right\vert}}
\newcommand{\ip}[2]{\ensuremath{\left\langle#1,#2\right\rangle}}
\newcommand{\p}{\ensuremath{\partial}}
\newcommand{\pr}{\mathcal{P}}

\newcommand{\pbar}{\ensuremath{\bar{\partial}}}
\newcommand{\db}{\overline\partial}
\newcommand{\D}{\mathbb{D}}
\newcommand{\B}{\mathbb{B}}
\newcommand{\Sp}{\mathbb{S}}
\newcommand{\T}{\mathbb{T}}
\newcommand{\R}{\mathbb{R}}
\newcommand{\Z}{\mathbb{Z}}
\newcommand{\C}{\mathbb{C}}
\newcommand{\N}{\mathbb{N}}
\newcommand{\Q}{\mathbb{Q}}
\newcommand{\mQ}{\mathcal{Q}}
\newcommand{\mS}{\mathcal{S}}
\newcommand{\scrH}{\mathcal{H}}
\newcommand{\scrL}{\mathcal{L}}
\newcommand{\td}{\widetilde\Delta}
\newcommand{\pw}{\text{PW}}
\newcommand{\esup}{\text{ess.sup}}
\newcommand{\Tn}{\mathcal{T}_n}
\newcommand{\Bn}{\mathbb{B}_n}
\newcommand{\rt}{\mathcal{O}}
\newcommand{\avg}[1]{\langle #1 \rangle}
\newcommand{\one}{\mathbbm{1}}
\newcommand{\eps}{\varepsilon}
\newcommand{\grad}{\nabla}

\newcommand{\La}{\langle }
\newcommand{\Ra}{\rangle }
\newcommand{\rk}{\operatorname{rk}}
\newcommand{\card}{\operatorname{card}}
\newcommand{\ran}{\operatorname{Ran}}
\newcommand{\osc}{\operatorname{OSC}}
\newcommand{\im}{\operatorname{Im}}
\newcommand{\re}{\operatorname{Re}}
\newcommand{\tr}{\operatorname{tr}}
\newcommand{\vf}{\varphi}
\newcommand{\f}[2]{\ensuremath{\frac{#1}{#2}}}

\newcommand{\kzp}{k_z^{(p,\alpha)}}
\newcommand{\klp}{k_{\lambda_i}^{(p,\alpha)}}
\newcommand{\TTp}{\mathcal{T}_p}
\newcommand{\m}[1]{\mathcal{#1}}
\newcommand{\md}{\mathcal{D}}
\newcommand{\qan}{\abs{Q}^{\alpha/n}}
\newcommand{\sbump}[2]{[[ #1,#2 ]]}
\newcommand{\mbump}[2]{\lceil #1,#2 \rceil}
\newcommand{\cbump}[2]{\lfloor #1,#2 \rfloor}

\newcommand{\class}{Math 629 Spring 26}
\newcommand{\term}{Fall 24}
\newcommand{\examnum}{Homework \hwnum}
\newcommand{\examdate}{\duedate}
\newcommand{\timelimit}{75 Minutes}
\newcommand{\vc}[3]{\langle #1,#2,#3\rangle}
\newcommand*{\vv}[1]{\vec{\mkern0mu#1}}
\newcommand{\bv}[1]{\boldsymbol{#1}}
\newcommand{\hide}[1]{}
\newcommand{\uvec}[1]{\boldsymbol{\hat{\textbf{#1}}}}
\newcommand{\vex}[1]{\boldsymbol{{\textbf{#1}}}}

\pagestyle{head}
\firstpageheader{}{}{}
\runningheader{\class}{\examnum\ - Page \thepage\ of \numpages}{\examdate}
\runningheadrule

\makeatletter
\renewcommand*\env@matrix[1][*\c@MaxMatrixCols c]{%
  \hskip -\arraycolsep
  \let\@ifnextchar\new@ifnextchar
  \array{#1}}
\makeatother

%\printanswers
\begin{document}

\noindent
\begin{tabular*}{\textwidth}{l @{\extracolsep{\fill}} r @{\extracolsep{6pt}} l}
\textbf{\class} & \textbf{Name: Mengxiang Jiang}\\
\end{tabular*}\\
\rule[2ex]{\textwidth}{2pt}
%

 

\newcommand{\duedate}{03/30/2026 at 11:59pm}
\newcommand\hwnum{10}

Turn this assginment in on gradescope. All work must be neat and legible. Turn in 
a ``final draft''. There should be nothing scratched out and your work should flow
generally from left to right and top to bottom.


\begin{questions}
\question 
In the text, it says that Newton was able to come up with the series 
for $\sqrt{1+x}$ using the "standard" root extraction algorithm. 
Find that algorithm and compute the series using Newton's Method. 
\\\\
Newton adapted the standard long square root algorithm (the pencil-and-paper method for extracting square roots digit-by-digit) to work on power series.
\\
To find $P = \sqrt{1+x}$, write $P = a_0 + a_1 x + a_2 x^2 + \cdots$ and build it term by term.

At each step you have a partial sum $p$ and a remainder $r = (1+x) - p^2$. The next term $\delta$ satisfies:
$$(p + \delta)^2 = p^2 + 2p\delta + \underbrace{\delta^2}_{\text{higher order}} \approx 1+x$$
So $\delta \approx \dfrac{r}{2p}$, and specifically $\delta = \dfrac{\text{leading term of } r}{2 \cdot \text{leading term of } p}$.

Then update: $r \leftarrow r - \bigl[(p+\delta)^2 - p^2\bigr] = r - 2p\delta - \delta^2$.
\\
The first few steps are as follows:\\
Step 1: $p = 1$ gives $p^2 = 1$. Remainder: $r = (1+x) - 1 = x$.

Step 2: $\delta = \dfrac{x}{2 \cdot 1} = \dfrac{x}{2}$. Update $p = 1 + \dfrac{x}{2}$.
$$r \leftarrow x - 2(1)\left(\tfrac{x}{2}\right) - \left(\tfrac{x}{2}\right)^2 = x - x - \frac{x^2}{4} = -\frac{x^2}{4}$$

Step 3: $\delta = \dfrac{-x^2/4}{2 \cdot 1} = -\dfrac{x^2}{8}$. Update $p = 1 + \dfrac{x}{2} - \dfrac{x^2}{8}$.

If we continue this process, we get the result:
$$\sqrt{1+x} = 1 + \frac{1}{2}x - \frac{1}{8}x^2 + \frac{1}{16}x^3 - \cdots$$

This matches the binomial series $(1+x)^{1/2}$ with the general term:
$$\binom{1/2}{n} x^n = \frac{\tfrac{1}{2}\left(\tfrac{1}{2}-1\right)\cdots\left(\tfrac{1}{2}-n+1\right)}{n!}x^n$$

\newpage 
\question 
Find the curvature of the ellipse $x^2 + 4y^2 = 1$ by using
Newton's procedure.
\\\\
Take fluxions of $x^2 + 4y^2 = 1$ with $\dot{x} = 1$:

$$2x + 8y\dot{y} = 0 \implies z = \dot{y} = -\frac{x}{4y}.$$

Take fluxions again (using $\dot{x} = 1$, $\dot{y} = z$):

$$2 + 8z^2 + 8y\dot{z} = 0 \implies \dot{z} = -\frac{1 + 4z^2}{4y}.$$

Substituting $z = -x/4y$, so $4z^2 = x^2/4y^2$:

$$1 + 4z^2 = \frac{4y^2 + x^2}{4y^2} = \frac{1}{4y^2}$$

where the last step uses $x^2 + 4y^2 = 1$. Thus $\dot{z} = -\dfrac{1}{16y^3}$.

Substituting $z^2 = x^2/16y^2$ and $x^2 = 1 - 4y^2$:

$$1 + z^2 = \frac{16y^2 + x^2}{16y^2} = \frac{12y^2 + 1}{16y^2}.$$

Newton's formula then gives:

$$DC = \frac{(1+z^2)^\frac{3}{2}}{\dot{z}} = \frac{\left(\dfrac{12y^2+1}{16y^2}\right)^{3/2}}{\dfrac{1}{16y^3}} = \frac{(12y^2+1)^{3/2}}{64y^3} \cdot 16y^3 = \frac{(12y^2+1)^{3/2}}{4}.$$

\newpage 
\question 
Derive the product rule using differentials. 
\\\\
Let $u$ and $v$ be two quantities depending on $x$. Consider their product $uv$ at $x$ and at $x + dx$:

$$d(uv) = (u + du)(v + dv) - uv.$$

Expanding:

$$d(uv) = uv + udv + vdu + dudv - uv.$$

The term $dudv$ is the product of two infinitesimals, and so is very small compared to the remaining terms. If we discard it:

$$d(uv) = udv + vdu.$$

\newpage
\question 
Using Leibniz's Method, find the tangent to $x^2+xy+y^2=7$ at the 
point $(1,2)$.
\\\\
Take the differential of both sides:

$$d(x^2) + d(xy) + d(y^2) = d(7).$$

Applying the rules for differentials (using the product rule for $d(xy)$):

$$2xdx + xdy + ydx + 2ydy = 0.$$

Collect terms in $dx$ and $dy$:

$$(2x + y)dx + (x + 2y)dy = 0.$$

Solving for the ratio $dy/dx$ (the slope of the tangent):

$$\frac{dy}{dx} = -\frac{2x+y}{x+2y}.$$

At the point $(1, 2)$:

$$\frac{dy}{dx}\bigg|_{(1,2)} = -\frac{2(1)+2}{1+2(2)} = -\frac{4}{5}.$$

The tangent line through $(1,2)$ with slope $-4/5$:

$$y - 2 = -\frac{4}{5}(x-1).$$
\end{questions}
\end{document}
